\documentclass[11pt,a4paper]{article}

% --- Engine: kompilovat pomocí XeLaTeX (kvůli českým znakům) ---
\usepackage{fontspec}
\usepackage{polyglossia}
\setdefaultlanguage{czech}
\setmainfont{Latin Modern Roman}
\setmonofont{Latin Modern Mono}

\usepackage[a4paper,margin=2.4cm]{geometry}
\usepackage{xcolor}
\usepackage{graphicx}
\usepackage{enumitem}
\usepackage{listings}
\usepackage{booktabs}
\usepackage{longtable}
\usepackage{titlesec}
\usepackage{fancyhdr}
\usepackage[hidelinks]{hyperref}

\definecolor{brand}{HTML}{2D6A4F}
\definecolor{gold}{HTML}{B8960C}
\definecolor{codebg}{HTML}{F4F6F5}
\definecolor{codeframe}{HTML}{D8DEDB}

\hypersetup{colorlinks=true, linkcolor=brand, urlcolor=brand}

\titleformat{\section}{\Large\bfseries\color{brand}}{\thesection}{0.6em}{}
\titleformat{\subsection}{\large\bfseries\color{brand!85!black}}{\thesubsection}{0.6em}{}

\pagestyle{fancy}
\fancyhf{}
\renewcommand{\headrulewidth}{0.4pt}
\lhead{\small\color{brand}PlaceSecrets}
\rhead{\small Příručka administrátora}
\cfoot{\small\thepage}

\lstset{
  basicstyle=\ttfamily\small,
  backgroundcolor=\color{codebg},
  frame=single,
  rulecolor=\color{codeframe},
  framesep=5pt,
  breaklines=true,
  breakatwhitespace=false,
  columns=fullflexible,
  keepspaces=true,
  showstringspaces=false,
  xleftmargin=4pt,
  literate={á}{{\'a}}1 {č}{{\v c}}1 {ď}{{\v d}}1 {é}{{\'e}}1 {ě}{{\v e}}1
           {í}{{\'i}}1 {ň}{{\v n}}1 {ó}{{\'o}}1 {ř}{{\v r}}1 {š}{{\v s}}1
           {ť}{{\v t}}1 {ú}{{\'u}}1 {ů}{{\r u}}1 {ý}{{\'y}}1 {ž}{{\v z}}1
}

\newcommand{\code}[1]{\texttt{#1}}
% Tolerovat delší řádky (kód/cesty) bez přetékání do okraje.
\setlength{\emergencystretch}{3em}
\sloppy

\begin{document}

% ---------------- Titulní strana ----------------
\begin{titlepage}
\centering
\vspace*{3cm}
{\Huge\bfseries\color{brand} PlaceSecrets\par}
\vspace{0.4cm}
{\LARGE Příručka administrátora\par}
\vspace{0.3cm}
{\large Instalace, správa, architektura a zabezpečení\par}
\vspace{2cm}
{\large POI AudioGuide \& Travel Log\par}
{\large Mobilní aplikace \textbullet{} Webový portál \textbullet{} Supabase backend\par}
\vfill
{\small Verze dokumentu k 14.\,6.\,2026\par}
\end{titlepage}

\tableofcontents
\newpage

% ================================================================
\section{Úvod}
Tento dokument je úplná příručka pro \textbf{administrátora a provozovatele}
systému PlaceSecrets. Po jeho přečtení má být jasné, \emph{co kde je}, \emph{jak
to nainstalovat a nasadit}, \emph{jak to spravovat} a \emph{jak je vyřešené
zabezpečení}.

Systém se skládá ze tří částí sdílejících jeden backend:
\begin{itemize}[nosep]
  \item \textbf{Mobilní aplikace} (Flutter, Android + iOS) -- pro koncové uživatele.
  \item \textbf{Webový portál} (React) -- veřejná část \emph{a} administrace.
  \item \textbf{Backend} (Supabase: PostgreSQL, Auth, Storage, Edge Functions).
\end{itemize}

\paragraph{Doprovodné dokumenty v repozitáři.} Tato příručka shrnuje a odkazuje na:
\code{app/BUILD.md} (build mobilu), \code{backend/DEPLOYMENT.md} (nasazení
backendu), \code{backend/PRODUCTION\_SETUP.md} (produkční data),
\code{RELEASE-CHECKLIST.md} (co musí udělat člověk vs.\ co je hotové v~kódu) a
\code{APP-GAP-ANALYSIS.md} (stav a~analýza).

% ================================================================
\section{Architektura systému}

\subsection{Přehled}
\begin{center}
\begin{tabular}{@{}ll@{}}
\toprule
\textbf{Vrstva} & \textbf{Technologie} \\
\midrule
Mobil & Flutter/Dart, Riverpod, GoRouter, Drift (SQLite), supabase\_flutter \\
      & flutter\_map (OSM), just\_audio + flutter\_tts, geofencing, Firebase FCM \\
Web   & Vite + React 18 + TypeScript + Tailwind CSS, react-i18next (11 jazyků) \\
Backend & Supabase: PostgreSQL + RLS, Auth, Storage, Edge Functions (Deno) \\
Platby & Stripe (mobil checkout/refund), Comgate (web), Google/Apple Pay, Revolut \\
\bottomrule
\end{tabular}
\end{center}

\subsection{Struktura repozitáře}
\begin{lstlisting}
place-secrets/
  app/        ... Flutter mobilni aplikace (balicek cz.tuma.placesecrets)
  web/        ... React webovy portal (verejna cast + /admin administrace)
  backend/    ... Supabase: migrace, edge funkce, seed
    supabase/migrations/   ... 64 SQL migraci (forward-only)
    supabase/functions/    ... 15 edge funkci (Deno/TypeScript)
  tex/        ... tento dokument
  .github/workflows/ci.yml ... CI (web, edge, mobile, e2e)
  APP-GAP-ANALYSIS.md, RELEASE-CHECKLIST.md ... stav a checklist
\end{lstlisting}

\subsection{Klíčové datové struktury (PostgreSQL)}
Hlavní tabulky: \code{cities}, \code{pois}, \code{countries}, \code{poi\_types},
\code{poi\_audio\_guides}, \code{poi\_ratings} (hodnocení), \code{user\_profiles},
\code{admin\_users} (administrátoři + PIN), \code{contributor\_profiles},
\code{acl\_permissions} (role \texttimes{} zdroj \texttimes{} R/W/D),
\code{payment\_config}, \code{payment\_receipts} (účtenky), \code{landing\_config}
(CMS úvodní stránky) a tabulky nároků: \code{user\_access},
\code{user\_poi\_purchases}, \code{user\_city\_premium\_purchases},
\code{user\_country\_purchases}, \code{free\_visit\_log}. Na všech je zapnuté
\textbf{Row Level Security (RLS)}.

\subsection{Edge funkce (15)}
\code{admin-api} (proxy všech admin operací), \code{verify-admin-pin},
\code{change-admin-pin}, \code{invite-contributor}, \code{confirm-payment},
\code{generate-receipt}, \code{create-payment-intent}, \code{create-poi-checkout},
\code{confirm-poi-payment}, \code{payment-webhook}, \code{stripe-refund},
\code{comgate-create}, \code{comgate-webhook}, \code{delete-user-data} (GDPR),
\code{migrate-anonymous-data}.

% ================================================================
\section{Instalace a první nasazení}
Nasazení má čtyři kroky: \textbf{backend}, \textbf{web}, \textbf{mobil} a
\textbf{doménové/účtové} úkony. Pořadí: nejdřív backend (na něm závisí web i mobil).

\subsection{Předpoklady (účty a nástroje)}
\begin{itemize}[nosep]
  \item Supabase projekt (produkční, případně i staging).
  \item Firebase projekt (push notifikace -- FCM).
  \item Apple Developer + Google Play Console (vydání aplikací).
  \item Platební brány: Stripe a/nebo Comgate (živé klíče).
  \item Hosting webu: Netlify nebo Vercel; doména \code{placesecrets.app}.
  \item Nástroje: \code{supabase} CLI, \code{flutter} SDK, \code{node} 22, \code{git}.
\end{itemize}

\subsection{Backend (Supabase)}
Podrobně v \code{backend/DEPLOYMENT.md}. Stručně (spouštět v \code{backend/}):
\begin{lstlisting}
npx supabase login
npx supabase link --project-ref <PROJECT_REF>
npx supabase db push          # aplikuje vsech 64 migraci

# tajne klice (viz tabulka v sekci Zabezpeceni):
npx supabase secrets set STRIPE_SECRET_KEY=... \
  COMGATE_MERCHANT=... COMGATE_SECRET=... \
  WEBHOOK_SECRET=$(openssl rand -hex 32) \
  RESEND_API_KEY=... APP_URL=https://placesecrets.app

# nasazeni vsech 15 edge funkci:
for fn in admin-api verify-admin-pin change-admin-pin \
    invite-contributor confirm-payment generate-receipt \
    create-payment-intent create-poi-checkout confirm-poi-payment \
    payment-webhook stripe-refund comgate-create comgate-webhook \
    delete-user-data migrate-anonymous-data; do
  npx supabase functions deploy "$fn"
done
\end{lstlisting}

\paragraph{Konfigurace Auth (dashboard).} V~\emph{Authentication}:
\begin{itemize}[nosep]
  \item \textbf{Redirect URLs}: přidat \code{placesecrets://login-callback} (OAuth z mobilu) a originy webu.
  \item \textbf{Providers}: zapnout Google a Apple (client id/secret).
  \item E-mailové šablony (potvrzení registrace, reset hesla).
\end{itemize}

\paragraph{Produkční data.} Podle \code{backend/PRODUCTION\_SETUP.md}: nasypat
\code{countries}, nastavit \code{payment\_config} (ceny, \code{free\_count}),
vytvořit prvního \code{superadmin}a, založit města a~místa (přes administraci).
Testovací seed (\code{seed.sql}) do produkce \textbf{nepatří}.

\subsection{Webový portál}
Deploy konfigurace je v~repozitáři hotová (\code{web/netlify.toml},
\code{web/vercel.json}, \code{web/public/\_redirects}). Na hostingu nastavte
proměnné prostředí (injektují se při buildu):
\begin{lstlisting}
VITE_SUPABASE_URL=https://<ref>.supabase.co
VITE_SUPABASE_ANON_KEY=<publikovatelny anon klic>   # chranen pomoci RLS
VITE_SENTRY_DSN=<volitelne, pro error tracking>
\end{lstlisting}
Build: \code{npm ci \&\& npm run build} (výstup \code{dist/}). SPA fallback i
bezpečnostní hlavičky (CSP, X-Frame-Options, \dots) jsou součástí deploy konfigurace.
Pokud chybí \code{VITE\_SUPABASE\_*}, aplikace se \textbf{záměrně} nespustí a~zobrazí
jasnou chybu (fail-fast v~\code{main.tsx}).

\subsection{Mobilní aplikace}
Podrobně v~\code{app/BUILD.md}. Člověkem dodávané části:
\begin{itemize}[nosep]
  \item \textbf{Firebase config}: \code{google-services.json} (Android) +
        \code{GoogleService-Info.plist} (iOS) + plugin \code{com.google.gms.google-services}.
        Bez nich appka \emph{funguje}, jen bez push notifikací (start je odolný).
  \item \textbf{Podpisový klíč}: vytvořit keystore a~vyplnit \code{android/key.properties}
        (mimo git). Bez něj se release podepíše debug klíčem (jen pro vývoj).
  \item \textbf{Ikona}: nahradit placeholder \code{assets/icons/app\_icon.png} a~spustit
        \code{dart run flutter\_\allowbreak launcher\_\allowbreak icons}.
  \item \textbf{Konfigurace}: vyplnit \code{dart\_defines/prod.json} (prod Supabase URL/klíč).
\end{itemize}
Build:
\begin{lstlisting}
flutter pub get
dart run build_runner build --delete-conflicting-outputs
flutter build appbundle --release --dart-define-from-file=dart_defines/prod.json
flutter build ipa       --release --dart-define-from-file=dart_defines/prod.json
\end{lstlisting}

\paragraph{Deep linky (doménové soubory).} Pro ověřené \code{https} odkazy hostujte
na doméně \code{/.well-known/\allowbreak assetlinks.json} (Android)
a~\code{apple-app-\allowbreak site-association} (iOS). Schéma \code{placesecrets://}
funguje i~bez nich (callbacky platby).

% ================================================================
\section{Webové administrační rozhraní}
Administrace je v~portálu pod \code{/admin/*} (přihlášení e-mailem + PIN).
Viditelnost položek řídí role přihlášeného administrátora (ACL).

\begin{longtable}{@{}p{4.6cm}p{9.8cm}@{}}
\toprule
\textbf{Sekce} & \textbf{Co se zde spravuje} \\
\midrule
\endhead
\code{/admin/dashboard} & Přehled, statistiky systému. \\
\code{/admin/cities} & Města (název, země, poloha, ceny, aktivace). \\
\code{/admin/cities/.../pois} & Místa (POI) ve městě: obsah, foto, audio, úroveň přístupu, ceny, slug. \\
\code{/admin/cities/.../pricing} & Cenotvorba města (bundle). \\
\code{/admin/poi-types} & Typy míst (12 typů) + lokalizace názvů. \\
\code{/admin/users} & Uživatelé: profil, stav (aktivní/deaktivovaný), GDPR export, granty přístupu. \\
\code{/admin/payments} & Platby a~účtenky, detail platby, refundace (Stripe). \\
\code{/admin/purchases/premium} & Přehled prémiových nákupů. \\
\code{/admin/contributors} & Přispěvatelé (autoři obsahu) a~jejich profily. \\
\code{/admin/hall-of-fame} & Žebříčky přispěvatelů (celkově / dle města / dle místa). \\
\code{/admin/moderation} & Moderace recenzí (skrýt/smazat nevhodné komentáře). \\
\code{/admin/landing} & CMS úvodní stránky: hero, Explore, Featured, How-it-works, CTA, patička -- per 11 jazyků. \\
\code{/admin/settings} & Systémová nastavení, platební konfigurace. \\
\code{/admin/acl} & Matice oprávnění: role \texttimes{} zdroj \texttimes{} čtení/zápis/mazání. \\
\code{/admin/my-pois} & Přispěvatel spravuje jen vlastní místa. \\
\code{/admin/profile} & Profil administrátora/přispěvatele, změna PIN. \\
\bottomrule
\end{longtable}

\paragraph{Role.} \code{superadmin} (vše), \code{editor} (POI + audio),
\code{viewer} (jen čtení), \code{contributor} (vlastní POI + profil),
\code{user} (čtení obsahu). Konkrétní práva jsou v~tabulce \code{acl\_permissions}
a~editovatelná na \code{/admin/acl}.

% ================================================================
\section{Zabezpečení}
Bezpečnostní model stojí na principu \emph{,,frontend nikdy nedrží silný klíč``}
a~\emph{,,databáze chrání sama sebe (RLS)``}.

\subsection{Autentizace a~autorizace}
\begin{itemize}[nosep]
  \item \textbf{Koncoví uživatelé}: anonymní session (okamžité použití) $\rightarrow$
        e-mail/heslo nebo OAuth (Google/Apple). Anonymní data se migrují při registraci.
  \item \textbf{Administrátoři}: e-mail + \textbf{PIN}. PIN je hashovaný (bcrypt),
        po~5 chybách 15min zámek (\code{verify-admin-pin}, atomické počítání pokusů).
        První přihlášení vynutí změnu PINu (\code{must\_change\_pin}).
  \item \textbf{ACL}: každá admin operace ověří právo (role \texttimes{} zdroj \texttimes{}
        R/W/D) v~\code{acl\_permissions}; \code{superadmin} má bypass.
\end{itemize}

\subsection{admin-api proxy -- žádný service\_role ve frontendu}
Frontend \textbf{nikdy} nepoužívá \code{service\_role} klíč. Všechny privilegované
operace jdou přes edge funkci \code{admin-api}, která: ověří JWT a~roli, validuje
názvy sloupců, blokuje citlivé sloupce (\code{pin\_hash} apod.), povoluje jen
whitelist tabulek a~RPC, vynucuje ACL a~loguje do \code{audit\_logs}.

\subsection{Row Level Security (RLS)}
Zapnuté na všech tabulkách. Uživatel vidí/edituje jen vlastní řádky
(\code{auth.uid()}); veřejná data (aktivní města/místa, \code{landing\_config})
mají \code{public read}; zápis přes \code{service\_role} (edge funkce) RLS obchází
řízeně.

\subsection{Storage (úložiště)}
Buckety: \code{poi-photos}, \code{audio-guides} (veřejné čtení), \code{avatars}
(veřejné čtení, zápis jen vlastník), \code{payment-receipts} (privátní -- servíruje
se přes krátkodobé \emph{signed URL}; navíc owner-read RLS policy jako
defense-in-depth), \code{premium-audio} (jen pro nakupující).

\subsection{Platby a~webhooky}
\begin{itemize}[nosep]
  \item \textbf{Stripe}: webhook ověřuje podpis (\code{constructEvent}); granty idempotentní.
  \item \textbf{Comgate}: webhook po kontrole sdíleného secretu \emph{re-ověřuje} transakci
        voláním Comgate status API (merchant+secret+transId) a~grantuje jen na ověřené
        \code{PAID}; refId/cena se berou z~ověřené odpovědi, ne z~těla callbacku.
  \item \textbf{Refundace}: jen administrátor s~ACL na \code{payments}; idempotentní (Stripe idempotencyKey).
\end{itemize}

\subsection{Hodnocení míst}
Hodnotit smí jen \textbf{přihlášený} (neanonymní) uživatel, \textbf{jednou na položku}
(UNIQUE \code{poi\_id,user\_id} + upsert). Anonymní zápis je odebrán na úrovni DB
(migrace 0453) i~aplikace.

\subsection{GDPR}
\begin{itemize}[nosep]
  \item \textbf{Export dat}: uživatel si stáhne svá data (web i~mobil).
  \item \textbf{Smazání účtu}: \code{delete-user-data} atomicky smaže data + storage + auth
        usera a~zapíše doklad do \code{audit\_logs} (\code{gdpr\_account\_deletion}).
  \item \textbf{Souhlasy}: časová razítka přijetí Podmínek/Soukromí při registraci;
        cookie lišta na webu (jen nezbytné úložiště, bez trackingu).
\end{itemize}

\subsection{Web -- bezpečnostní hlavičky}
Deploy konfigurace nastavuje \code{Content-Security-Policy}, \code{X-Frame-Options},
\code{X-Content-Type-Options}, \code{Referrer-Policy} a~\code{Permissions-Policy}.

\subsection{Správa tajných klíčů}
\begin{longtable}{@{}p{5.2cm}p{9.2cm}@{}}
\toprule
\textbf{Kde} & \textbf{Co} \\
\midrule
\endhead
Supabase secrets & \code{STRIPE\_SECRET\_KEY}, \code{STRIPE\_POI\_WEBHOOK\_SECRET},
  \code{COMGATE\_MERCHANT/SECRET}, \code{WEBHOOK\_SECRET}, \code{RESEND\_API\_KEY} -- jen na serveru. \\
Web (build env) & \code{VITE\_SUPABASE\_URL}, \code{VITE\_SUPABASE\_ANON\_KEY} (publikovatelný), \code{VITE\_SENTRY\_DSN}. \\
Mobil (dart-define) & \code{SUPABASE\_URL}, \code{SUPABASE\_ANON\_KEY} v~\code{dart\_defines/prod.json}. \\
Mobil (keystore) & \code{android/key.properties} -- mimo git. \\
\bottomrule
\end{longtable}
\textbf{Nikdy} necommitovat reálné klíče; v~repu jsou jen \code{*.example} šablony.

% ================================================================
\section{Údržba a~provoz}

\subsection{CI/CD}
GitHub Actions (\code{.github/workflows/ci.yml}) na každý push/PR spouští:
\begin{itemize}[nosep]
  \item \textbf{web}: lint, \code{tsc}, build, \code{vitest} (293 testů), \code{npm audit}.
  \item \textbf{edge-functions}: \code{vitest} (72 testů).
  \item \textbf{mobile}: \code{flutter analyze} + \code{flutter test} (233 testů) + codegen.
  \item \textbf{e2e}: Playwright (5 testů veřejných toků v~prohlížeči).
\end{itemize}

\subsection{Monitoring}
Web: Sentry (volitelné, aktivní jen když je nastaveno \code{VITE\_SENTRY\_DSN}).
Mobil: doporučeno Firebase Crashlytics (vyžaduje Firebase config). Edge chyby
jsou v~logu Supabase.

\subsection{Migrace databáze}
Konvence: \textbf{forward-only}, nepřečíslovávat. Nová změna = nový soubor
\code{NNNN\_popis.sql} (idempotentní -- \code{IF NOT EXISTS} apod.). Aplikace:
\code{npx supabase db push}. Lokálně po editaci: \code{npx supabase db reset}.

\subsection{Zálohy a~obnova}
Spoléhat na Supabase \emph{Point-in-Time Recovery} (dle plánu projektu). Před
velkými zásahy ruční záloha (export). Doporučen oddělený \textbf{staging} projekt
pro testy nasazení.

\subsection{Časté administrátorské úkoly}
\begin{itemize}[nosep]
  \item \textbf{Přidat město / místo}: \code{/admin/cities} $\rightarrow$ POI editor (foto, audio, ceny, slug).
  \item \textbf{Udělit přístup uživateli}: \code{/admin/users/:id} (grant město/POI/země).
  \item \textbf{Moderace}: \code{/admin/moderation} (skrýt/smazat komentář).
  \item \textbf{Změnit ceny / free limit}: \code{/admin/settings} (\code{payment\_config}).
  \item \textbf{Upravit úvodní stránku}: \code{/admin/landing} (per jazyk).
  \item \textbf{Přidat administrátora}: pozvánka (\code{invite-contributor}) / \code{/admin/users}.
\end{itemize}

% ================================================================
\section{Řešení problémů}
\begin{longtable}{@{}p{6cm}p{8.4cm}@{}}
\toprule
\textbf{Příznak} & \textbf{Příčina / řešení} \\
\midrule
\endhead
Web ukazuje ,,Configuration error`` & Chybí \code{VITE\_SUPABASE\_*} v~build env hostingu. \\
Přímý odkaz na \code{/admin/...} vrací 404 & Chybí SPA fallback -- ověřit \code{netlify.toml}/\code{vercel.json}. \\
Platba projde, ale přístup se neobjeví & Webhook nedoběhl/nedorazil; ověřit endpoint a~secrety brány; stránka úspěchu sama poolluje nárok. \\
Mobil spadne hned po startu & Chybí Firebase config; ověřit \code{google-services.json}/plist (start je jinak odolný). \\
Push notifikace nedorazí & Není FCM token na serveru / chybí Firebase config / oprávnění notifikací. \\
Edge funkce vrací 503 & Chybí příslušný secret (např. Comgate/Stripe) -- \code{supabase secrets set}. \\
Admin se nepřihlásí (zámek) & 15min zámek po 5 chybných PINech; počkat nebo reset přes superadmina. \\
\bottomrule
\end{longtable}

% ================================================================
\section{Odkazy}
\begin{itemize}[nosep]
  \item \code{app/BUILD.md} -- build mobilu, Firebase, keystore, deep linky, OAuth.
  \item \code{backend/DEPLOYMENT.md} -- nasazení backendu (migrace, secrets, funkce).
  \item \code{backend/PRODUCTION\_SETUP.md} -- produkční data (countries, admin, ceny).
  \item \code{RELEASE-CHECKLIST.md} -- co musí udělat člověk vs.\ co je hotové v~kódu.
  \item \code{APP-GAP-ANALYSIS.md} -- stav, analýza, ověření.
\end{itemize}

\vfill
\noindent\rule{\linewidth}{0.4pt}\\
{\small Tento dokument popisuje stav systému k~datu na titulní straně. Runtime chování
závislé na účtech a~zařízeních (platby, push, deep linky, store buildy) je nutné
ověřit v~cílovém prostředí.}

\end{document}
